\section{Analysis}\label{sec:theory}
%% One page. Two analyses only: (i) the mechanism, what each
%% supervision channel optimizes and how the advantage weight allocates
%% the budget; (ii) the deployment coverage guarantee. Everything else
%% lives in Appendix~\ref{app:analysis}. No em dashes, no colons in
%% prose, no magic numbers.

\paragraph{What each supervision channel optimizes.}
The two data channels of Section~\ref{sec:distill} induce updates with
qualitatively different expected drift. Maximum-likelihood training on
teacher trajectories follows
$-\nabla_\theta D_{\mathrm{KL}}(\pi_\phi\|\pi_\theta)$ per query, so
its expected update moves the student toward the teacher's behavior
distribution as a whole, and nothing in the objective distinguishes
the movement required by the task from the movement that merely
reproduces source-specific regularities. For trajectories drawn from
the student's own policy the score-function expectation vanishes,
$\mathbb{E}_{\tau\sim\pi_\theta}[\nabla_\theta\log\pi_\theta(\tau\mid
x)]=0$, and weighting the same samples by the verifier yields the
policy-gradient identity
\begin{equation}
\nabla_\theta J_x(\theta)
=\mathbb{E}_{\tau\sim\pi_\theta}\!\left[V(x,\tau)\,
\nabla_\theta\log\pi_\theta(\tau\mid x)\right],
\label{eq:pg-identity}
\end{equation}
so the only systematic drift of the verifier-filtered channel is the
ascent direction of task success. Both identities are standard; the
proposition below states what the advantage weight adds on top of
them.

\begin{proposition}[Deficit-proportional allocation with a safe fixed
point]\label{prop:advantage}
Let $\hat p(x)$ be the empirical unguided success rate of task $x$ and
let the positive objective weight each verified trajectory of $x$ by
$a(x)=(1-\hat p(x))_{+}$. Then in expectation over collection,
(i) the update direction contributed by task $x$ is the
policy-gradient ascent direction of Eq.~\eqref{eq:pg-identity} scaled
by $a(x)$, so the training budget concentrates on tasks in proportion
to their measured deficit;
(ii) if $\hat p(x)=1$ for every support task, every weight is zero,
the update vanishes, and the returned policy is the base student; and
(iii) as $\hat p(x)\to 0$ uniformly, the objective reduces to the
unweighted verifier-filtered objective.
\end{proposition}

The proof is immediate from linearity of the objective in $a(x)$ and
is given in Appendix~\ref{app:analysis}. Point (ii) is the property
the experiments exercise most; on a benchmark where the base student
already matches its attainable ceiling, the method's correct output is
the base policy itself, and the pipeline emits a zero-update decision
trace rather than a degraded checkpoint. Point (iii) states graceful
degradation toward the pure self-improvement regime when the student
fails everywhere, in which case guided collection supplies the
verified trajectories that make the objective non-vacuous.

\paragraph{Deployment coverage.}
For exchangeable in-task queries, the split-conformal threshold
calibrated on $\mathcal{S}_c$ yields the standard finite-sample
marginal coverage guarantee \citep{vovk2005algorithmic,
angelopoulos2023conformal}. The guarantee concerns acceptance of
in-task queries only; rejection power out of task is an empirical
quantity, and we report it as such.
